% Seccion: Convolución directa y por FFT
% Responsable: R4. Extension: 1,0 pagina(s).
% Solo R4 edita este archivo.
%
% Que debe contener (segun el enunciado y la pauta):
%   - Ecuaciones de la convolución lineal y del teorema de la convolución.
%   - Tiempos de la directa vs. FFT según el largo del filtro (convolucion_tiempos.png).
%   - Convolución circular sin relleno de ceros (convolucion_circular.png).

\section{Convolución directa y por FFT}
\label{sec:convolucion}

La convolución lineal de $x$ (largo $N$) con $h$ (largo $M$) es
\begin{equation}
  y[n] = \sum_{k} x[k]\, h[n-k], \qquad n = 0, \dots, N+M-2,
  \label{eq:conv}
\end{equation}
y entrega $N+M-1$ muestras. \texttt{convolucion\_directa} la calcula con dos ciclos, y el ciclo de
adentro recorre $h$ para que el costo sea $O(NM)$.

\texttt{convolucion\_fft} usa el teorema de la convolución, que dice que convolucionar en el tiempo
es lo mismo que multiplicar en frecuencia. Si $\tilde{x}$ y $\tilde{h}$ son $x$ y $h$ completadas
con ceros hasta largo $L = N+M-1$,
\begin{equation}
  y = \mathrm{IDFT}\big(\mathrm{DFT}(\tilde{x}) \cdot \mathrm{DFT}(\tilde{h})\big),
  \label{eq:teorema-conv}
\end{equation}
donde el producto es punto a punto y usamos la \texttt{dft} y la \texttt{idft} de la
Sección~\ref{sec:dft}. Al final nos quedamos con la parte real, porque la parte imaginaria que
queda es casi cero ($10^{-16}$) y sale solo por el redondeo. En las pruebas comparamos las dos
versiones con \texttt{np.convolve} y entre ellas, y revisamos que un impulso en la posición $d$
mueve la señal $d$ muestras y que $x*h = h*x$.

\paragraph{Por qué rellenar con ceros.} La DFT de $L$ puntos supone que la señal se repite cada $L$
muestras, por eso al multiplicar DFTs sale la convolución \emph{circular}. Si no se rellena, las
últimas $M-1$ muestras no caben y se suman al comienzo. En la Figura~\ref{fig:conv}a usamos un
pulso cerca del final de $x$ ($N=40$) y una exponencial de $M=15$ coeficientes para que la cola
sea grande. Se ve que la cola de la lineal (muestras 40 a 53) aparece igual en las muestras 0 a 13
de la circular, donde la lineal vale cero.

\paragraph{Tiempos.} Dejamos $N=500$ fijo y probamos $M$ entre 4 y 512. Cada caso se midió 3 veces
y usamos el menor tiempo, para que no influyeran otros programas del computador. En escala log-log
la directa es una recta de pendiente 1, o sea que su tiempo crece igual que $M$, como dice $O(NM)$
(Figura~\ref{fig:conv}b). La versión por DFT casi no cambia, porque su costo depende de
$L = N+M-1$ y ese largo cambia poco mientras $M$ sea mucho menor que $N$. En nuestra prueba las
curvas se cruzan cerca de $M = 128$. Hay que tener cuidado con ese número, porque nuestra
\texttt{dft} es $O(L^2)$ pero usa operaciones rápidas de NumPy, mientras que la directa usa ciclos
de Python, que son más lentos. Por eso lo que conviene comparar son las pendientes. Con una FFT de
verdad, que es $O(L\log L)$, el cruce pasaría con filtros mucho más cortos.

\begin{figure}[H]
  \centering
  \begin{minipage}{0.55\linewidth}\centering
    \includegraphics[width=\linewidth]{convolucion_circular.png}\\(a)
  \end{minipage}\hfill
  \begin{minipage}{0.42\linewidth}\centering
    \includegraphics[width=\linewidth]{convolucion_tiempos.png}\\(b)
  \end{minipage}
  \caption{(a) Convolución lineal y circular sin relleno. La franja roja es la cola que se pasó
  al comienzo. (b) Tiempo de las dos versiones según el largo de $h$.}
  \label{fig:conv}
\end{figure}

Para el caso C agregamos \texttt{convolucion\_separable\_2d}, que pasa un filtro por cada fila y
otro por cada columna de la imagen. El resultado es el mismo que usar el filtro 2D
$h_\text{cols}\,h_\text{filas}^{T}$ (lo comparamos con \texttt{scipy.signal.convolve2d}), pero se
hacen $2M$ operaciones por píxel en vez de $M^2$. La usamos en el prefiltro de la
Sección~\ref{sec:aliasing}.
